% Resume: Arumalla Mohit Krishna
% Compile with pdfLaTeX.

\documentclass[letterpaper,10pt]{article}

\usepackage[margin=0.5in]{geometry}
\usepackage[T1]{fontenc}
\usepackage{lmodern}
\usepackage{microtype}
\usepackage{titlesec}
\usepackage{enumitem}
\usepackage[hidelinks]{hyperref}
\usepackage{tabularx}
\input{glyphtounicode}

\pdfgentounicode=1
\pagestyle{empty}
\urlstyle{same}
\raggedright
\setlength{\parindent}{0pt}
\setlength{\tabcolsep}{0pt}

\titleformat{\section}
  {\large\scshape}{}{0em}{}[\titlerule]

\titlespacing*{\section}{0pt}{6pt}{3pt}

\newcommand{\entry}[2]{%
  \noindent
  \begin{tabularx}{\linewidth}
    {@{}>{\raggedright\arraybackslash}X@{\hspace{6pt}}r@{}}
    \textbf{#1} & #2
  \end{tabularx}\par
}

\newcommand{\bullets}{
  \begin{itemize}[
    leftmargin=1.2em,
    topsep=2pt,
    itemsep=1pt,
    parsep=0pt,
    partopsep=0pt
  ]
}

\newcommand{\finishbullets}{
  \end{itemize}
}

\begin{document}

%----------HEADING----------
\begin{center}
    {\LARGE\bfseries Arumalla Mohit Krishna}\\[3pt]
    +1 (631) 352-8643 $|$
    \href{mailto:mohitkrishna.arumalla@stonybrook.edu}
    {mohitkrishna.arumalla@stonybrook.edu} $|$
    \href{https://linkedin.com/in/mohit-krishna-arumalla}
    {LinkedIn}
\end{center}
\vspace{-8pt}

%-----------EDUCATION-----------
\section{Education}

\entry
  {\textbf{Stony Brook University}}
  {Aug 2025 -- Dec 2026}

M.S. in Computer Science \hfill New York, USA\\
\textit{
Coursework: Natural Language Processing,
Introduction to Computer Vision, Vision Language Modeling
}
\vspace{3pt}

\entry
  {\textbf{Birla Institute of Technology and Science, Pilani}}
  {Aug 2019 -- 2023}

B.E. in Electronics \& Communication; GPA: 8.02/10 \hfill India
%-----------PROFESSIONAL SUMMARY-----------
\section{Professional Summary}

Software engineer with 2+ years of experience building distributed services,
semantic search, and real-time data pipelines, pursuing an M.S. in Computer
Science at \textbf{Stony Brook University}. Hands-on experience in personalized
retrieval, GPU model serving, and hallucination detection, with ongoing work
in self-supervised vision and parameter-efficient adaptation. Skilled in
\textbf{Python, Go, PyTorch, AI-assisted development tools, and cloud infrastructure}.
Proven ability to \textbf{build production-level code, collaborate cross-functionally, and iterate rapidly}.

%-----------TECHNICAL SKILLS-----------
\section{Technical Skills}

\textbf{Languages:}
Python, Go, C++, SQL, Java, TypeScript, JavaScript\\

\textbf{ML \& Retrieval:}
PyTorch, TensorFlow/Keras, scikit-learn, NumPy, pandas,
Transformers, mBART, BGE, Vision Transformers,
Self-Supervised Learning, LoRA, Visual Prompt Tuning,
Text Embeddings, Vector Search\\

\textbf{Data \& Systems:}
gRPC, Kafka, Apache Flink, Flink SQL, Debezium/CDC,
PostgreSQL, Redis, Elasticsearch, Airflow, BigQuery\\

\textbf{Infrastructure:}
NVIDIA Triton, AWS, Docker, Kubernetes, Linux, Git,
CI/CD, Istio/Envoy, New Relic\\

\textbf{AI Tools:}
Claude Code, Codex, Cursor, AI-assisted development environments

%-----------EXPERIENCE-----------
\section{Experience}

\entry
  {\textbf{Software Engineer II $|$ Infra.Market, Bangalore}}
  {Apr 2024 -- Aug 2025}

\bullets

  \item Built \textbf{production-level semantic product search} using Amazon Titan text embeddings
  and Elasticsearch vector search, enabling natural-language product
  discovery through chat interfaces.

  \item Engineered a \textbf{low-latency CDC pipeline} with Flink, Kafka, Debezium,
  PostgreSQL, RocksDB, and Elasticsearch; improved synchronization latency
  to approximately 500 ms with exactly-once guarantees.

  \item Deployed Istio/Envoy across 15+ microservices and integrated
  New Relic distributed tracing for production monitoring, ensuring \textbf{high system reliability and security}.

\finishbullets

\entry
  {\textbf{Software Engineer I $|$ Infra.Market, Chennai}}
  {Jul 2023 -- Mar 2024}

\bullets

  \item Owned Search and Location microservices; built gRPC services
  handling 500K+ daily requests for product discovery and geolocation,
  with a focus on scalability and maintainability.

  \item Fixed PostgreSQL connection leaks using Gorm ORM, reducing
  connection-pool exhaustion by 85\%; validated with concurrent
  load tests at 10K requests/second.

  \item Built Airflow pipelines for BigQuery-to-Elasticsearch migration
  with query optimization, dynamic filtering, and automated monitoring,
  supporting \textbf{robust and automated data workflows}.

\finishbullets

%-----------PROJECTS-----------
\section{Projects}

\entry
  {GPU-Optimized Personalized Product Search}
  {Aug 2026 -- Oct 2026}

\textit{Python, PyTorch, BGE, NVIDIA Triton, Go, gRPC}

\bullets

  \item Developed two-tower retrieval with
  \texttt{BAAI/bge-small-en-v1.5} query/product embeddings
  (384 dimensions) and a trainable user tower combining query
  embeddings with interaction history.

  \item Built an inference pipeline with precomputed product embeddings,
  approximate nearest-neighbor retrieval, and a Go/gRPC gateway;
  integrating Triton for GPU query-encoder serving.

  \item Designed batching and concurrent-instance benchmarks for
  steady/bursty traffic, measuring p95/p99 latency, throughput,
  GPU utilization, and cost per successful search; evaluated
  Recall@K/NDCG@K while testing overload recovery.

\finishbullets

\entry
  {Attention-Based Hallucination Detection in NMT}
  {Apr 2026 -- May 2026}

\textit{Python, PyTorch, mBART, Transformers, scikit-learn}

\bullets

  \item Built an English-German hallucination-detection pipeline using
  teacher-forced mBART attention features from WMT22/WMT23 quality
  estimation data.

  \item Extracted cross-/self-attention features across 12 decoder
  layers and 16 heads, including entropy, max attention mass,
  variance, top-k mass, and self-to-cross grounding ratios.

  \item Trained token-, sentence-, and head-level classifiers,
  achieving 0.712 F1 and 0.695 ROC-AUC on 50K WMT examples;
  proposed hallucination-aware reranking using localization,
  risk scores, and attention signals.

\finishbullets

\entry
  {Prompt/LoRA for Self-Supervised Learning}
  {Ongoing}

\textit{Python, PyTorch, Vision Transformers}

\bullets

  \item Investigating domain-specific visual prompting and LoRA
  in a self-supervised ViT framework to learn transferable
  representations across image domains.

  \item Studying backbone performance without prompts/LoRA at
  inference; designing comparisons with retained or combined modules
  to evaluate classification performance and parameter efficiency
  under domain shift.

\finishbullets

\end{document}